\section{Discussion}
\label{sec:discussion}

The principal lesson of the archived workflow is that transfer in this problem is governed less by generic expressive power than by how efficiently a descriptor maps local transition-metal bonding onto the small data regime. Bond-order-potential moments were designed precisely to summarize the local electronic structure of metallic environments, and in Fe--Mo TCP phases that physical bias matters. Even though the stored internal test table gives a slightly smaller numerical RMSE for ACE with kernel ridge regression, the bond-specific BOP representation reaches almost the same internal error from a far smaller and more interpretable starting feature space, and it is the representation that the validation and thermodynamics notebooks carry forward to the complex phases. In other words, the BOP model is not attractive because it always wins the in-sample scoreboard by a narrow margin, but because it packages the chemistry most relevant for transfer.

The importance of CNAV reinforces this interpretation. Conventional unit-cell averaging would erase the distinction between, for example, a CN12 site and a CN16 site even though those environments occupy very different roles in a TCP lattice. The archived comparisons show that once that distinction is preserved, the error falls sharply for the descriptors that already contain local-bonding information. The roughly 24\% BOP improvement and the roughly 40\% SOAP improvement under CNAV are therefore not accidental numerical gains. They show that a large part of the learning task is the recovery of crystallographic site identity, and that a representation respecting that identity transfers more faithfully from the simpler prototypes to the chemically richer $R$, $P$, $M$, and $\delta$ structures.

This explains why a simple-to-complex learning strategy can work at all. The complex TCP phases do not invent entirely new local motifs; rather, they recombine familiar coordination polyhedra and bond topologies across a larger number of inequivalent sites. If the descriptor resolves those local motifs and if the model is trained on a sufficiently broad range of simple structures, then the extrapolation challenge is softened into a controlled interpolation in the space of local environments. The complex-phase validation notebook operationalizes exactly this idea: it does not rely on a fully exhaustive DFT map of the large-cell structures, but instead checks whether a model trained on simpler environments retains accuracy when those environments are embedded into more elaborate sublattice arrangements.

The thermodynamic step is especially important because it tests more than pointwise energy accuracy. Bragg--Williams free-energy minimization amplifies systematic energetic biases into errors in predicted sublattice occupancies. The fact that the archived workflow proceeds to an experimentally anchored occupancy comparison for the $R$ phase implies that the energy model is stable enough to support finite-temperature inference rather than only static ranking. This is a stronger standard than test-set regression alone. It also suggests that the learned Hamiltonian contains enough physically meaningful composition dependence to interact sensibly with a compound-energy-formalism model. At the same time, the present repository only supports an end-member-based mean-field treatment, so any agreement with experiment should be interpreted as evidence that the dominant site preference has been captured, not as proof that short-range-order effects are negligible.

The present workflow nevertheless has clear limitations. First, the archived training and validation calculations are explicitly nonmagnetic in the principal workflow, and magnetic contributions could matter for some Fe-rich structures or for phase-boundary shifts. Second, the Bragg--Williams treatment neglects short-range order beyond the mean field, which is acceptable for a first pass but cannot capture all correlation effects in multisublattice phases. Third, several of the most interesting structures lie close in energy, so residual errors of a few tens of meV/atom may still reorder near-degenerate variants. Finally, the checked-out repository preserves the validation logic and figures more clearly than the raw serialized parity metrics, which complicates the extraction of a fully tabulated external-error matrix from the notebook JSON alone.

These limitations do not undercut the broader significance of the approach. Direct DFT enumeration of all occupation patterns in large-cell TCP phases would require a prohibitive number of relaxations, whereas the present workflow distills the problem into a physically structured representation learning task. Because the domain knowledge enters through generic ideas---Vegard scaling for chemistry, coordination-resolved averaging for crystallography, and moment-based bonding descriptors for local electronic structure---the same recipe should be extendable to other binary transition-metal systems and, with suitable care, to ternaries. We nevertheless treat that broader applicability as an outlook rather than as a demonstrated result of the present repository. In that sense, the Fe--Mo study is best viewed as a prototype for data-efficient surrogate modeling of complex intermetallic free-energy landscapes rather than as a one-off regression exercise.
